\scanpage{127}
% archive foliation: 126
 où on trouve \deleted{un groupe algébrique \uline{réel} $\mathrm{GP}(1)_{\mathbb R}$}\inserted{le groupe des points réels $\mathrm{GP}(1,\mathbb R)$} --- qui en fait est une forme réelle du $\mathrm{GP}(1)_{\mathbb C}$, \deleted{\unsure{donnant un stabilisateur}}\inserted{\unsure{d'ailleurs il se trouve}} qu'le \deleted{choix d'une}\inserted{\unsure{correspond}} \deleted{forme} biunivoque entre

$\alpha)$ choix d'une \deleted{forme réelle}\inserted{descente du corps de base de $\mathbb C$ à $\mathbb R$} dans le groupe alg. \deleted{$\unsure{\mathrm{GP}_{\mathbb C}}$} $\operatorname{Aut}_{\mathbb C}P$

$\beta)$ choix d'une descente du corps de base de $\mathbb C$ à $\mathbb R$ dans l'algèbre de Lie complexe
\[
\operatorname{Lie}\operatorname{Aut}_{\mathbb C}P=\mathrm H^0(P,\mathcal T_{P/\mathbb C})\simeq\mathrm H^0(P,\Omega^{-1}_{P/\mathbb C})
\]
$\gamma)$ choix d'un domaine homogène de type $3^\circ$ (i.e. non algébrique) dans $P$, à polarité près (la polaire de $D$ est $\complement\overline D$).

\margin{$\delta)$ choix d'un \unsure{géodésique} \unsure{n.|non} ponctuel ($=$ cercle) dans $P$}
La relation entre $\alpha)$ et $\beta)$ claire, celle entre \unsure{$\delta)$|$\alpha)$} et $\gamma)$ est la suivante : À une \deleted{\unsure{déforme}} projective réelle $\inserted{P_{\mathbb R}}$, on associe l'ens. de ses pts réels, qui est un ``cercle'' dans $P_{\mathbb R}(\mathbb C)=P(\mathbb C)$ $(P=P_{\mathbb R}\otimes_{\mathbb R}\mathbb C)$, et on prend les deux composantes connexes de $P_{\mathbb R}(\mathbb C)-P_{\mathbb R}(\mathbb R)\ldots$ Notons que l'inversion, par un cercle $P_{\mathbb R}(\mathbb R)$

\uline{Scholie} On a un foncteur $\inserted{P}$ de la catégorie des courbes analytiques non sing. connexes et simplement connexes $\inserted{X}$ vers la catégorie des droites projectives complexes, et une immersion fonctorielle
\[X\lhook\joinrel\xrightarrow{i_X}P(X)(\mathbb C)\]
faisant de $X$ un domaine \unsure{strict.} hom. de $P(X)$.
Un tel foncteur est unique (à iso unique près respectant $i:\operatorname{id}\longrightarrow P(-)(\mathbb C)$).

Soit maintenant $X$ une courbe analytique non singulière. Pour tt pt $x\in X$, soit $\widetilde X(x)$ le rev. universel de $X$ en $x$. C'est un espace analytique simplement connexe, et quand $x$ varie, on reçoit un système local (foncteur du groupoïde fondamental $\Pi(X)$ de $X$) à valeurs les $\widetilde X(x)$ (considérés comme objets de la cat des \deleted{esp.}\inserted{courbes} analytiques connexes et simplement connexes, \unsure{munis par pt donné}). Le système des projections
\[\widetilde X(x)\longrightarrow X=X(x)\]
est un hom. de systèmes locaux d'espaces analytiques. On lui associe donc toujours un hom. d'espaces analytiques sur $X$
\diagram{\begin{tikzcd}
\widetilde X\arrow[r]\arrow[dr]&X\times X\arrow[d,"\operatorname{pr}_2"]\\
&X
\end{tikzcd}}
munis de connexions intégrables (au sens strict : il y a transport parallèle\ldots)
L'hom. est compatible avec les connexions, et transforme la section marquée $\sigma:x\mapsto\widetilde x$ de $\widetilde X$ en la section diagonale de $X$ (d'où on conclut que $\sigma$ est analytique).

Cette construction marche pour tt espace analytique. Mais utilisant le fait que $\deleted{X}$ est une courbe non sing., donc les $\widetilde X(x)$ sont des courbes non sing. conn. et simplement connexes, on trouve que le système local des $\widetilde X(x)$ donne naissance au système local des $P(\widetilde X(x))(\mathbb C)$, d'où un \uline{fibré projectif canonique}, $\mathcal P$, sur
